La literatura previa en análisis predictivo de Churn se ha enfocado primariamente en la maximización del rendimiento algorítmico, favoreciendo ensambles de árboles (por ejemplo, Random Forest o Gradient Boosting) sobre modelos lineales tradicionales dada su destreza en esquemas de datos tabulares con relaciones no lineales.

A pesar de su éxito estadístico, múltiples teóricos coinciden en una barrera comercial: la opacidad algorítmica. Para mitigarlo, investigadores en \textit{eXplainable Artificial Intelligence} (XAI) propusieron implementaciones \textit{post-hoc} agnósticas al modelo, siendo LIME y SHAP (SHapley Additive exPlanations) los de mayor adopción empírica por su robustez apoyada en la teoría de juegos cooperativos \cite{lundberg2017unified}. 

Simultáneamente, la perspectiva académica referente al Costo de Clasificación Errante destaca que en el entorno empresarial no todos los errores tienen el mismo peso. Penalizar falsos positivos (campañas de retención gastadas en vano) y falsos negativos (el ingreso de vida del cliente o \textit{LTV} perdido) demanda enfoques \textit{Profit-driven} \cite{verbeke2012new}. A pesar de estos avances aislados, existe una notoria brecha en la literatura contemporánea respecto a estudios aplicados que operen las predicciones y sus desgloses causales XAI dentro de sistemas interactivos de Inteligencia de Negocios (BI) para optimización financiera en tiempo real.
